\section{Results: does Fisher--Rao help ABF?}
\label{sec:results}

\subsection{Headline comparison}
\Cref{tab:main_integrated} reports the four methods, each at its best
configuration under the \emph{integrated}-$\Ltwo(F)$ rule (convergence speed;
see \cref{sec:selection}). Reading the practical comparison
(ABF only versus ABF+FR estimated):
the estimated-target FR reduces the median integrated free-energy error from
$\AbfIntegF$ to $\EstIntegF$, a $\EstIntegImpPct\%$ reduction, while the
\emph{final} free-energy error improves only slightly
(from $\AbfFinalF$ to $\EstIntFinalF$). The birth--death is extremely gentle:
the median per-event replacement fraction is $\EstIntEvt$, i.e.\ of order one
particle in $10^{5}$ per FR application. \Cref{fig:convergence} shows the same
story over time: the FR curves (estimated, uniform, oracle) detach below the
ABF-only curve in the early and middle phase of the $\Ltwo(F)$ panel and the
ensemble re-equilibrates faster, with the bands (inter-quartile range over
seeds) overlapping by the end of the budget.

\begin{table}[t]
\centering
\small
\caption{Main comparison, configurations selected by smallest median \emph{integrated} $L^2(F)$ (convergence speed). Lower is better except $P(\text{beats ABF})$.}
\label{tab:main_integrated}
\begin{tabular}{lrrrrr}
\toprule
Method & $\int_0^T\!\|\hat F_t-F_{\rm ref}\|\,dt$ & $\|\hat F_T-F_{\rm ref}\|$ & $\|\hat F'_T-F'_{\rm ref}\|$ & $P(\text{beats ABF})$ & FR event frac. \\
\midrule
ABF only & 15.91 & 0.0137 & 0.0718 & -- & $0$ \\
FR estimated & 14.08 & 0.0126 & 0.0771 & 0.60 & $1.6\times10^{-5}$ \\
FR uniform & 14.46 & 0.0125 & 0.0769 & 0.80 & $5.3\times10^{-6}$ \\
FR oracle* & 13.27 & 0.0113 & 0.0602 & 1.00 & $1.7\times10^{-5}$ \\
\bottomrule
\end{tabular}
\par\smallskip\footnotesize $^{*}$\,The oracle target uses the reference free energy $F_{\rm ref}$ (the unknown being computed); it is a diagnostic positive control, \emph{not} a usable algorithm. Per-config medians over $5$ seeds; one configuration per method, selected by the smallest median integrated $L^2(F)$.
\end{table}


\begin{figure}[t]
  \centering
  \includegraphics[width=0.95\textwidth]{fig02_convergence_curves.png}
  \caption{Convergence of the selected (integrated-$\Ltwo(F)$ rule)
    configuration for each method: median (line) and inter-quartile range
    (band) over \NSeeds{} seeds. The free-energy and mean-force errors
    (top, log scale) fall fastest for the oracle diagnostic and faster for the
    estimated/uniform FR than for ABF only in the early-to-middle phase;
    barrier crossings (bottom left) are essentially indistinguishable, and the
    FR event fraction (bottom right) is of order $10^{-5}$ throughout ---
    FR is acting as a weak correction, not a violent resampling. The oracle
    curve is a diagnostic control (\cref{rem:oracle}).}
  \label{fig:convergence}
\end{figure}

\subsection{Final profiles}
\Cref{fig:profiles} overlays the final mean-force, free-energy and
reaction-coordinate-marginal profiles. All practical methods recover
$\Fpref$ and $\Fref$ well; the visible differences between ABF only and
ABF+FR estimated at the final budget are small, consistent with the modest
final-error improvements in \cref{tab:main_integrated}. The $x$-marginal
(panel c) is somewhat flatter under FR, as intended, but --- as we emphasise
in \cref{sec:metrics} --- marginal flatness is a diagnostic, not the goal.

\begin{figure}[t]
  \centering
  \includegraphics[width=0.97\textwidth]{fig03_final_profiles.png}
  \caption{Final profiles (mean over seeds) for the selected configurations.
    (a) Mean force $\Fphat(x)$ and (b) centred free energy $\Fhat(x)$ versus the
    reference (dashed); (c) reaction-coordinate marginal $\phat(x)$, with the
    unbiased $\pref$ (dotted) for reference. Practical methods recover the
    reference well; FR mildly flattens $\phat$.}
  \label{fig:profiles}
\end{figure}

\subsection{Target-type ablation}
\Cref{fig:ablation} shows the distribution, over the \NFRconfigs{}
configurations of each target type, of the seed-median errors. Three honest
observations follow. First, for the \emph{integrated} error all three FR target
types have medians below the ABF-only baseline (dashed line) --- FR typically
accelerates convergence --- but the spread is wide and some configurations are
\emph{worse} than ABF, so FR does not help unconditionally. Second, for the
\emph{final} $\Ltwo(F)$ the estimated and uniform medians sit essentially on the
ABF baseline: the practical final-budget gain is marginal. Third, and most
telling, on the final \emph{mean-force} error the estimated and uniform medians
lie slightly \emph{above} the ABF baseline (a small degradation at the typical
configuration), whereas only the oracle control is consistently below it. In
other words, flattening the marginal with an imperfect target buys early speed
but not final mean-force accuracy; only an accurate target direction (the
oracle) improves everything at once.

\begin{figure}[t]
  \centering
  \includegraphics[width=0.97\textwidth]{fig04_target_ablation.png}
  \caption{Target-type ablation: each box is the distribution over the
    \NFRconfigs{} configurations of a target type (each point a seed-median);
    the dashed line is the ABF-only median. FR lowers the \emph{integrated}
    error (left) but the practical targets barely change the \emph{final}
    $\Ltwo(F)$ (middle) and slightly raise the final mean-force error (right) at
    the median configuration. Only the oracle diagnostic improves all three.}
  \label{fig:ablation}
\end{figure}

\subsection{Two selection rules give different ``best'' hyperparameters}
\label{sec:selection}
Because integrated and final errors measure different things, they need not be
optimised by the same configuration, and conflating them would be misleading.
We therefore report two rules explicitly:
\begin{itemize}
  \item selecting by \emph{integrated} $\Ltwo(F)$
        (\cref{tab:main_integrated}, the main text) favours \emph{fast
        convergence}; for the estimated target it selects an
        \emph{aggressive} configuration ($\gamma=\EstIntGamma$, burn-in
        $\EstIntBurnin$);
  \item selecting by \emph{final} $\Ltwo(F)$
        (\cref{tab:main_final}, supplement) favours \emph{final-budget
        accuracy}; for the estimated target it selects a \emph{gentle, delayed}
        configuration ($\gamma=\EstFinGamma$, burn-in $\EstFinBurnin$), which
        improves the final free-energy error by $\EstFinFinalImpPct\%$ and the
        final mean-force error by $\EstFinFinalFpImpPct\%$ over ABF only, and
        beats ABF on $\EstFinProb\%$ of matched seeds.
\end{itemize}
This dependence on the selection rule is itself a finding: an aggressive FR rate
accelerates the integrated error but does not (here) improve final accuracy,
while a gentle, late FR is the safer choice for final accuracy. We emphasise the
integrated rule in the main text because the clearest, most robust effect of FR
in this example is on convergence \emph{speed}; we are explicit that the
final-budget improvements are modest.

\begin{table}[t]
\centering
\small
\caption{Supplement: configurations selected by smallest median \emph{final} $L^2(F)$ (final-budget accuracy).}
\label{tab:main_final}
\begin{tabular}{lrrrrr}
\toprule
Method & $\int_0^T\!\|\hat F_t-F_{\rm ref}\|\,dt$ & $\|\hat F_T-F_{\rm ref}\|$ & $\|\hat F'_T-F'_{\rm ref}\|$ & $P(\text{beats ABF})$ & FR event frac. \\
\midrule
ABF only & 15.91 & 0.0137 & 0.0718 & -- & $0$ \\
FR estimated & 15.25 & 0.0120 & 0.0673 & 0.80 & $1.1\times10^{-6}$ \\
FR uniform & 14.97 & 0.0118 & 0.0680 & 0.60 & $7.3\times10^{-6}$ \\
FR oracle* & 13.27 & 0.0113 & 0.0602 & 1.00 & $1.7\times10^{-5}$ \\
\bottomrule
\end{tabular}
\par\smallskip\footnotesize $^{*}$\,Oracle is a diagnostic control (uses $F_{\rm ref}$), not a deployable method. Configurations here are selected by smallest median \emph{final} $L^2(F)$; cf.\ Table~\ref{tab:main_integrated}, which need not pick the same hyperparameters.
\end{table}


\subsection{The uniform target is a competitive naive ablation}
In this comparatively easy example the \emph{uniform} target is competitive with
the estimated target (integrated reduction $\UniIntegImpPct\%$; best final
$\Ltwo(F)=\UniFinalF$). We do not read this as evidence that uniform flattening
is a general solution: it is a naive ablation that happens to work here because,
under a near-converged bias $\Bt\approx F$, the canonical biased marginal is
already close to uniform, so ``flatten $\phat$'' and ``match the biased
marginal'' nearly coincide. In harder regimes (\cref{sec:discussion}) the two
targets would diverge and the uniform target would be expected to misdirect the
birth--death.

\subsection{The oracle control indicates headroom, not a method}
The oracle diagnostic (\cref{rem:oracle}) achieves the lowest errors of all ---
integrated reduction $\OracIntegImpPct\%$, final reduction $\OracFinalImpPct\%$,
and it beats ABF on $\OracProb\%$ of matched seeds. \emph{This is not a claim
that one should use the oracle target}; the oracle is unusable because it
requires $\Fref$. The correct reading is comparative: the estimated target
captures $\EstIntegImpPct\%$ of the integrated-error reduction while the oracle
reaches $\OracIntegImpPct\%$, so a substantial part of the potential benefit of
the FR mechanism is left unrealised by the current online target estimate. The
oracle--estimated gap thus localises the practical bottleneck in the
\emph{target estimation}, not in the birth--death mechanism itself.

\subsection{Answer to Q1}
\label{sec:q1}
On this example: \emph{yes, but modestly, and chiefly as a convergence
accelerator}. Practical estimated-target FR improves the time-integrated
free-energy error appreciably ($\EstIntegImpPct\%$) and gives a smaller
final-budget improvement; the gain is real but not a dramatic rescue, and it is
not unconditional --- poorly chosen FR configurations match or underperform ABF.
FR is therefore best described, in this regime, as a \emph{convergence
accelerator and stabiliser} rather than a method that qualitatively changes what
ABF can compute. The oracle control shows that more is possible \emph{if} the
target can be estimated better.
